% !TeX root = ../动手学强化学习笔记.tex
% 由主文档通过 \input 载入；请编译主文档。
\setcounter{chapter}{3} % 保留原书章号。
\chapter{马尔可夫决策过程}
\label{ch:rl-mdp}

赌博机只要求估计动作的即时收益，而一般决策必须同时考虑动作如何改变将来能够获得的信息与奖励。马尔可夫决策过程（Markov decision process，MDP）把这种影响写成可递推的概率模型。它描述问题，不指定求解算法。

\section{马尔可夫性质与状态建模}

令完整历史为 $H_t=(S_0,A_0,R_1,\ldots,S_t)$。受控马尔可夫性质要求在所有有定义的条件事件上，
\[
\Pr(S_{t+1}=s',R_{t+1}=r\mid H_t,A_t=a)
=p(s',r\mid S_t,a).
\]
右侧与更早历史及时间 $t$ 无关，分别体现状态信息充分与动力学平稳。这里不是说相邻状态独立，而是说给定当前状态与动作后，早期历史不再增加下一步预测信息。

\begin{example}[位置为何可能不是状态]
小车采用 $X_{t+1}=X_t+V_t$、$V_{t+1}=V_t+A_t$ 的离散运动模型。两条历史可能在时刻 $t$ 到达同一位置却具有不同速度；即使采用同一动作，下一位置仍不同。因此单独 $X_t$ 一般不满足所需马尔可夫性质，$(X_t,V_t)$ 则包含该模型下一步演化所需的信息。
\end{example}

没有动作时，状态转移构成马尔可夫过程；再附加奖励成为马尔可夫奖励过程（Markov reward process，MRP）。固定一个平稳策略后，MDP 也会诱导一个 MRP。观测是实际获取的信息，状态是用于预测未来的充分表示；二者相同是一种模型假设。

\section{有限 MDP 的组成与策略类别}

\begin{definition}[有限折扣 MDP]
一个有限折扣 MDP 由有限状态集 $\mathcal S$、每个状态上非空有限动作集 $\mathcal A(s)$、平稳联合转移奖励核 $p(s',r\mid s,a)$、初始状态分布 $\mu$ 和折扣因子 $0\leq\gamma<1$ 构成。奖励几乎处处有界。为便于书写，离散奖励情况下要求 $p\geq0$ 且 $\sum_{s',r}p(s',r\mid s,a)=1$；连续奖励时把对 $r$ 的求和换成积分。
\end{definition}

边缘转移和期望奖励分别为
\[
P(s'\mid s,a)=\sum_r p(s',r\mid s,a),\qquad
r(s,a)=\sum_{s',r}r\,p(s',r\mid s,a).
\]
在 $P(s'\mid s,a)>0$ 时，还可定义 $r(s,a,s')=\mathbb E[R_{t+1}\mid s,a,s']$。联合核允许奖励与后继状态相关，不能一般地拆成奖励分布和转移分布的乘积。小写 $r$ 在求和中为奖励取值，在 $r(s,a)$ 中为期望奖励函数，应由自变量区分。

平稳马尔可夫策略 $\pi(a\mid s)$ 不显式依赖时间；非平稳策略 $\pi_t(a\mid s)$ 可以随剩余时域改变。历史依赖策略写为 $\pi_t(a\mid H_t)$。确定性策略每次指定一个动作，随机策略给出动作分布。有限折扣 MDP 的控制可由平稳确定性策略达到最优，但有限时域任务的最优动作通常还依赖时间。

\section{回合任务、持续任务与终止约定}

若任务在停止时刻 $T$ 终止，回报为 $G_t=\sum_{k=0}^{T-t-1}\gamma^kR_{t+k+1}$。真实终止后不再产生任务收益，可以添加奖励为零的吸收状态，将后续价值设为零。固定有限时域下允许 $\gamma=1$；随机终止且无折扣时，仅有 $T<\infty$ 几乎必然还不够保证回报可积，奖励有界及 $\mathbb E[T]<\infty$ 是一组充分条件。

采样截断表示程序暂时停止收集轨迹，不一定意味着任务终止。如果持续任务的采样窗口在时刻 $h$ 结束，则前 $h$ 步奖励之后仍有未来价值；将其误设为零会改变训练目标。反之，若“最多运行 $h$ 步”就是任务规则，时限是实际终止条件，且剩余时间可能需要并入状态。

持续任务还可使用平均奖励 $g_\pi=\lim_{T\to\infty}T^{-1}\mathbb E_\pi[\sum_{t=0}^{T-1}R_{t+1}]$。该极限的存在性和对初始状态的依赖需另作假设；例如固定策略诱导的有限不可约马尔可夫链具有唯一平稳分布，期望平均奖励可由该分布加权计算。它不是把无限折扣和中的 $\gamma$ 直接设为一所得的有限数。

\section{从模型到环境实现}

一个网格环境可以把可通行格子作为状态，上下左右作为动作；撞墙保持原位，移动通常奖励 $-1$，到达目标后终止。如果动作以概率 $0.8$ 按意图移动、其余概率滑向侧面，就得到随机转移。奖励应明确对应“离开格子”“进入格子”还是整次转移，否则同一轨迹的回报可能被不同程序计算成不同数值。

读取全部 $P,r$、能够从任意指定状态动作生成样本，以及只能沿当前轨迹与环境交互，是三种不同的信息条件。前者支持精确动态规划，中者常称拥有生成模型访问，后者最接近在线交互。环境程序的重置分布、非法动作处理和随机数机制都属于实验定义；算法不应读取任务没有开放的隐状态或真实参数。

\section{习题}
\begin{enumerate}
\item 某设备状态为正常或故障，正常时继续运行以概率 $0.1$ 故障，收益为 $3$；检修花费 $1$ 并使下一状态正常。补充故障时的合法动作和收益，写出一套完整 MDP，并区分自己的建模选择与题设条件。
\item 给出随机终止时间 $T$ 几乎必然有限但 $\mathbb E[T]=\infty$ 的分布。若每步奖励为 $1$ 且无折扣，说明其价值为何不有限。
\item 在固定五步的任务中，同一位置在剩余一步与剩余四步时的最优动作能否不同？构造一个例子，并说明怎样扩充状态。
\end{enumerate}





